\section{Conclusion}
\label{sec:conclusion}

This paper investigates whether an offline reinforcement learning agent equipped with
a CVaR reward function and macroeconomic state variables can produce equity portfolios
with superior risk-adjusted performance out-of-sample relative to classical benchmarks
on the Brazilian stock exchange B3.
We formulate the portfolio rebalancing problem as a Markov Decision Process, construct
a 225-dimensional state vector that encodes asset-level return dynamics together with
SELIC regime indicators, the policy rate level and the bank lending spread, and train a
Conservative Q-Learning (CQL) agent on an offline dataset generated by the
Minimum-Variance Portfolio behavioral policy.
Evaluated on 720 out-of-sample trading days shared with the benchmarks, the full
proposed model is dominated by an equal-weight portfolio on every metric reported:
Sharpe (0.4058 against 0.6076), Sortino, CVaR$_{0.95}$, maximum drawdown and
Calmar.
Risk Parity dominates it as well.
The $1/N$ puzzle documented by \citet{DeMiguel2007} holds in this setting.

The value of the study lies in the ablation, which localizes the failure rather
than reporting it in aggregate.

\begin{enumerate}
  \item \textbf{Conservatism is the component that works.}
        Raising the CQL conservative weight from approximately zero to its tuned
        value is the single largest effect in the study, worth $+0.29$ in Sharpe,
        and it leaves tail loss unchanged.
        Bounding an offline policy to actions the behavioral data supports is
        doing real work on this dataset.
        We note that the ablation does not compare online with offline training:
        every configuration is offline, and no online agent was trained.

  \item \textbf{The CVaR reward, as specified, is not a reward.}
        Replacing the realised-return signal with the CVaR objective costs $-0.24$
        in Sharpe and, contrary to its purpose, worsens realised tail loss, which
        rises monotonically across the two configurations that use it.
        The cause is structural: the objective scores candidate weights against the
        same twenty-day window the state already encodes, making it a deterministic
        function of the observation.
        An agent maximising it receives no feedback about the consequences of its
        actions and converges to a rolling short-window tail optimizer fitted to a
        single tail observation.
        This is the finding we consider most useful to other practitioners, because
        the error is invisible at the level of the equations and appears only when
        the reward and the state are read side by side.

  \item \textbf{Regime conditioning is unresolved, and the evidence available
        points the wrong way.}
        The macro block recovers $+0.20$ in Sharpe but not the tail deterioration,
        and its regime-conditional pattern inverts the motivating hypothesis: the
        agent trails the naive benchmark most widely during tightening, and leads
        only while rates fall.
        That signature is what a persistently defensive allocation produces without
        learning anything about monetary policy.
        Because A3 and A4 share the defective reward, the comparison between them
        cannot settle the question either way.
\end{enumerate}

A fourth result is operational rather than statistical.
The evaluation applies the policy daily, giving measured turnover of 38 to 88 times
portfolio value per year against roughly one for the benchmarks.
At ten basis points per round trip this is 4--9\% annually, against agent returns
of 2.6--8.0\%.
No comparison in this paper is informative about implementable performance, and we
report none as such.

The order of future work follows from the above, and the first two items are
prerequisites rather than extensions.
First, the reward must be re-specified so that it scores an action against returns
the agent has not observed, and A3 and A4 re-run; until then the paper's second and
third questions remain open rather than answered.
Second, the ablation must be repeated across random seeds and reported with
dispersion: every number here comes from one seed and one window, and the effects
discussed are within the range seed variation alone can produce.
Third, turnover must be constrained in the objective or charged in the evaluation,
since the current protocol compares a daily-rebalancing agent with quarterly
benchmarks at zero cost.
Beyond those, extending the asset universe to include fixed-income instruments,
currency futures, and real-estate investment trusts (FIIs) would test whether the
framework generalizes across asset classes.
Second, incorporating realistic transaction costs and market-impact models would
provide a more accurate picture of the net-of-cost performance improvement
available to institutional investors.
Third, replacing the rule-based SELIC regime classifier with a real-time macro
nowcasting model or a hidden Markov Model could reduce the regime-labeling lag
and further improve performance during transition periods.
Finally, exploring alternative offline RL algorithms---such as Decision Transformer
\citep{Chen2021} or IQL \citep{Kostrikov2021}---within the same framework would
broaden the evidence base for offline RL as a tool for financial portfolio
optimization in emerging markets.
